% Part: first-order-logic
% Chapter: axiomatic-deduction
% Section: proof-theoretic-notions

\documentclass[../../../include/open-logic-section]{subfiles}

\begin{document}

\iftag{FOL}
      {\olfileid{fol}{axd}{ptn}}
      {\olfileid{pl}{axd}{ptn}}

\olsection{مفاهیم برهان‌نظری}

\begin{explain}
همان‌گونه که شماری از مفاهیم مهم معنایی
(\iftag{FOL}{اعتبار}{همان‌گویی}، استلزام، ارضاپذیری) را تعریف کرده‌ایم،
اکنون \emph{مفاهیم برهان‌نظری} متناظر را تعریف می‌کنیم. این مفاهیم با
توسل به ارضای !!{sentence}s در
\iftag{FOL}{!!{structure}s}{!!{valuation}s} تعریف نمی‌شوند، بلکه با
توسل به !!{derivability} یا !!{nonderivability} برخی فرمول‌ها تعریف
می‌شوند. اینکه این مفاهیم بر هم منطبق‌اند کشفی مهم بود. انطباق آن‌ها
مضمون \emph{قضیه‌های صحت} و \emph{تمامیت} است.
\end{explain}

\begin{defn}[!!^{derivability}]
!!^a{formula}~$!A$ از $\Gamma$~\emph{!!{derivable}} است، و
$\Gamma \Proves !A$ نوشته می‌شود، اگر !!a{derivation} از~$\Gamma$ وجود
داشته باشد که به~$!A$ پایان یابد.
\end{defn}

\begin{defn}[قضیه‌ها]
!!^a{formula}~$!A$ یک \emph{قضیه} است اگر !!a{derivation} برای $!A$ از
مجموعهٔ تهی وجود داشته باشد. اگر $!A$ قضیه باشد، $\Proves !A$، و اگر
نباشد، $\Proves/ !A$ می‌نویسیم.
\end{defn}

\begin{defn}[سازگاری]
مجموعهٔ $\Gamma$ از !!{formula}s \emph{سازگار} است هرگاه و تنها هرگاه
$\Gamma\Proves/ \lfalse$؛ و در غیر این صورت \emph{ناسازگار} است.
\end{defn}

\begin{prop}[بازتابی‌بودن]
\ollabel{prop:reflexivity}
اگر $!A \in \Gamma$، آنگاه $\Gamma \Proves !A$.
\end{prop}

\begin{proof}
  !!{formula}~$!A$ به‌تنهایی !!a{derivation} برای~$!A$ از~$\Gamma$ است.
\end{proof}

\begin{prop}[یکنواختی]
\ollabel{prop:monotonicity}
اگر $\Gamma \subseteq \Delta$ و $\Gamma \Proves !A$، آنگاه
$\Delta \Proves !A$.
\end{prop}

\begin{proof}
هر !!{derivation} برای $!A$ از $\Gamma$، !!a{derivation} برای $!A$ از
~$\Delta$ نیز هست.
\end{proof}

\begin{prop}[تعدی]
\ollabel{prop:transitivity}
اگر $\Gamma \Proves !A$ و $\{!A\} \cup \Delta \Proves !B$، آنگاه
$\Gamma \cup \Delta \Proves !B$.
\end{prop}

\begin{proof}
  فرض کنید $\{!A\} \cup \Delta \Proves !B$. آنگاه
  !!a{derivation}~$!B_1$، \dots، $!B_l = !B$ از~$\{!A\} \cup
  \Delta$ وجود دارد. برخی گام‌های آن !!{derivation} به‌سبب قاعده‌ای
  درست‌اند که به سطر پیشین~$!B_i = !A$ ارجاع می‌دهد. بنا بر فرض،
  !!a{derivation} برای~$!A$ از~$\Gamma$ وجود دارد؛ یعنی
  !!a{derivation}~$!A_1$، \dots، $!A_k = !A$ که در آن هر $!A_i$ یا
  اصل موضوع است، یا !!a{element}~$\Gamma$، یا به‌وسیلهٔ قاعدهٔ استنتاج
  درست است. اکنون دنبالهٔ زیر را در نظر بگیرید:
  \[
  !A_1, \dots, !A_k = !A, !B_1, \dots, !B_l = !B.
  \]
  این !!{derivation} درستی برای~$!B$ از $\Gamma \cup \Delta$ است، زیرا
  اکنون هر $!B_i = !A$ با همان قاعده‌ای توجیه می‌شود که
  ~$!A_k = !A$ را توجیه می‌کند.
\end{proof}

توجه کنید این حکم به‌ویژه بدان معناست که اگر $\Gamma \Proves !A$ و
$!A \Proves !B$، آنگاه $\Gamma \Proves !B$. همچنین از آن نتیجه می‌شود
که اگر $!A_1, \dots, !A_n \Proves !B$ و برای هر~$i$،
$\Gamma \Proves !A_i$، آنگاه $\Gamma \Proves !B$.

\begin{prop}
\ollabel{prop:incons}
$\Gamma$ ناسازگار است هرگاه و تنها هرگاه برای هر~$!A$،
$\Gamma \Proves !A$.
\end{prop}

\begin{proof}
تمرین.
\end{proof}

\tagprob{FOL}
\begin{prob}
\olref[fol][axd][ptn]{prop:incons} را اثبات کنید.
\end{prob}
\tagendprob

\tagprob{notFOL}
\begin{prob}
\olref[pl][axd][ptn]{prop:incons} را اثبات کنید.
\end{prob}
\tagendprob

\begin{prop}[فشردگی]
\ollabel{prop:proves-compact}
  \begin{enumerate}
  \item اگر $\Gamma \Proves !A$، آنگاه زیرمجموعه‌ای متناهی
    $\Gamma_0 \subseteq \Gamma$ وجود دارد چنان‌که
    $\Gamma_0 \Proves !A$.
  \item اگر هر زیرمجموعهٔ متناهی~$\Gamma$ سازگار باشد، آنگاه $\Gamma$
    سازگار است.
  \end{enumerate}
\end{prop}

\begin{proof}
  \begin{enumerate}
    \item اگر $\Gamma \Proves !A$، آنگاه دنبالهٔ متناهی
      $!A_1$، \dots،~$!A_n$ از !!{formula}s وجود دارد چنان‌که
      $!A \ident !A_n$ و هر $!A_i$ یا اصل موضوعی منطقی است، یا
      !!a{element}~$\Gamma$، یا به‌وسیلهٔ یک قاعدهٔ استنتاج از
      !!{formula}s پیشین به دست می‌آید. $\Gamma_0$ را متشکل از همان
      $!A_i$هایی بگیرید که در~$\Gamma$ هستند. آنگاه !!{derivation} نیز
      !!a{derivation} از~$\Gamma_0$ است و بنابراین
      $\Gamma_0 \Proves !A$.
    \item این حکم عکس نقیض~(۱) برای حالت خاص
      $!A \ident \lfalse$ است.
\end{enumerate}
\end{proof}

\end{document}
